\section{Lecture IV: Hitchin fibration} 

During the second half of 1980s, Hitchin introduced the famous integrable system, the moduli space of Higgs bundles, in his study of gauge theory. Around the same time, Kazhdan and Lusztig introduced affine Springer fibers as natural analogs of Springer fibers. For more than 15 years these two objects stayed unrelated until B-C.Ng\^o saw a connection between the two. Ng\^o's fundamental insight can be summarized as saying that Hitchin fibers are global analogs of affine Springer fibers, while affine Springer fibers are local versions of Hitchin fibers. Here ``global'' refers to objects involving a global function field, or an algebraic curve, rather than just a local function field, or a formal disk.  This observation, along with ingenious technical work, allowed Ng\^o to prove the Fundamental Lemma for orbital integrals conjectured by Langlands and Shelstad. We will review Hitchin's integrable system in a slightly more general setting, and make precise its connection to affine Springer fibers.




\subsection{The Hitchin moduli stack}

\subsubsection{The setting} We are back to the setting in \S\ref{ss:notation}. In addition, we fix an algebraic curve $X$ over $k$ (assumed algebraically closed) which is smooth, projective and connected.


\subsubsection{The moduli stack of bundles}
There is a moduli stack $\Bun_{n}$ classifying vector bundles of rank $n$ over $X$.  For any $k$-algebra $R$, $\Bun_{n}(R)$ is the groupoid of rank $n$ vector bundles (locally free coherent sheaves) on $X_{R}:=X\times_{\Spec k}\Spec R$. The stack $\Bun_{n}$ is algebraic, see \cite[Th 4.6.2.1]{LMB}. Moreover, it is smooth and locally of finite type over $k$.

\subsubsection{$G$-torsors}
Recall a {\em (right) $G$-torsor} over $X$ is a scheme $\calE$ over $X$ with a fiberwise action of $G$, such that locally for the \'etale topology of $X$, $\calE$ becomes the $G\times X$ and the $G$-action becomes the right translation action of $G$ on the first factor. 



For general reductive $G$, We have the moduli stack $\Bun_{G}$ of $G$-torsors over $X$. For a $k$-algebra $R$, the $R$-points of $\Bun_{G}$ is the groupoid of $G$-torsors over $X_{R}$. Then $\Bun_{G}$ is also a smooth algebraic stack locally of finite type over $k$.

\subsubsection{Associated bundles}\label{sss:ass bun} Let $(V,\rho)$ be a $k$-representation of $G$. Let $\calE$ be a $G$-torsor over $X$. Then there is a vector bundle $\rho(\calE)$ of $X$ whose total space is
\begin{equation*}
\textup{Tot}(\rho(\calE))=G\bs (\calE\times V)
\end{equation*}
where $G$ acts on $(e,v)\in\calE\times V$ by $g\cdot(e,v)=(eg^{-1}, \rho(g)v)$. The vector bundle $\rho(\calE)$ is said to be associated to $\calE$ and $\rho$.

When $G=\GL_{n}$, there is an equivalence of groupoids
\begin{equation}\label{GLn vec}
\{ \mbox{vector bundles $\calV$ of rank $n$ over $X$}\}\cong\{\mbox{$\GL_{n}$-torsors $\calE$ over $X$}\}. 
\end{equation}
The direction $\calV\mapsto\calE$ sends a vector bundle $\calV$ to the $\GL_{n}$-torsor of framings of $\calV$, namely take $\calE=\un{\Isom}_{X}(\calO^{n}_{X},\calV)$, with the natural action of $\GL_{n}$ on the trivial bundle $\calO^{n}_{X}$. The other direction $\calE\mapsto\calV$ sends a $\GL_{n}$-torsor $\calE$ over $X$ to the vector bundle $\St(\calE)$ associated to $\calE$ and the standard representation $\St$ of $\GL_{n}$. The equivalence \eqref{GLn vec} gives a canonical isomorphism of stacks $\Bun_{n}\cong\Bun_{\GL_{n}}$.


\subsubsection{}\label{sss:Bun Sp} For other classical groups $G$, $G$-torsors have more explicit descriptions in terms of vector bundles. For example, when $G=\SL_{n}$, a $G$-torsor over $X$ amounts to the same thing as a pair $(\calV,\iota)$ where $\calV$ is a vector bundle over $X$ of rank $n$ and $\iota$ is an isomorphism of line bundles $\wedge^{n}\calV\cong\calO_{X}$.


When $G=\Sp_{2n}$, the groupoid of $G$-torsors $\calE$ on $X$ is equivalent to the groupoid of pairs $(\calV,\omega)$ where $\calV$ is a vector bundle of rank $2n$ over $X$ and $\omega:\wedge^{2}\calV\to\calO_{X}$ is an $\calO_{X}$-linear map of coherent sheaves that gives a symplectic form on geometric fibers. The map in one direction sends a $\Sp_{2n}$-torsor $\calE$ to the pair $(\calV,\omega)$, where $\calV=\St(\calE)$ is the vector bundle associated to $\calE$ and the standard representation $\St$ of $\Sp_{2n}$, and the symplectic form $\omega$ on $\calV$ comes from the canonical map of $\Sp_{2n}$-representations $\wedge^{2}(\St)\to \one$ (where $\one$ is the trivial representation). 



\subsubsection{Higgs bundles} Fix a line bundle $\calL$ over $X$. An {\em $\calL$-twisted Higgs bundle of rank $n$} over $X$  is a pair $(\calV,\varphi)$ where $\calV$ is a vector bundle over $X$ of rank $n$, and $\varphi:\calV\to\calV\otimes\calL$ is an $\calO_{X}$-linear map. The endomorphism $\varphi$ is called a {\em Higgs field} on $\calV$.

There is a moduli stack $\calM_{n,\calL}$ classifying $\calL$-twisted Higgs bundles of rank $n$ over $X$. The morphism $\calM_{n,\calL}\to \Bun_{n}$ forgetting the Higgs field is representable. Therefore $\calM_{n,\calL}$ is also an algebraic stack over $k$.


\subsubsection{$G$-Higgs bundles} An {\em $\calL$-twisted $G$-Higgs bundle} is a pair $(\calE,\varphi)$ where $\calE$ is a $G$-torsor over $X$ and $\varphi$ is a global section of the vector bundle $\Ad(\calE)\otimes\calL$ over $X$. Here, $\Ad(\calE)$ is the vector bundle associated to $\calE$ and the adjoint representation $(\frg,\Ad)$ of $G$, in the sense of \S\ref{sss:ass bun}. We call $\varphi$ an {\em $\calL$-twisted Higgs field} on $\calE$.

When $G=\GL_{n}$, the notion of $\calL$-twisted $G$-Higgs bundle is equivalent to that of an $\calL$-twisted Higgs bundle of rank $n$. In fact, to each $\calL$-twisted $G$-Higgs bundle $(\calE,\varphi)$, we get a Higgs bundle $(\calV=\St(\calE),\phi)$, where $\phi:\calV\to\calV\otimes\calL$ viewed as a global section of $\underline{\End}(\calV)\otimes\calL$ corresponds to $\varphi$ under the canonical isomorphism $\underline{\End}(\calV)\cong\Ad(\calE)$.

We also have the moduli stack $\calM_{G,\calL}$ of $\calL$-twisted Higgs $G$-torsors over $X$. The $R$-points of $\calM_{G,\calL}$ is the groupoid of $\calL_{R}$-twisted $G$-Higgs bundles on $X_{R}$, where $\calL_{R}$ denotes the pullback of $\calL$ to $X_{R}$. The forgetful morphism $\calM_{G,\calL}\to \Bun_{G}$ is representable, hence $\calM_{G,\calL}$ is an algebraic stack over $k$. When $G=\GL_{n}$, we have a canonical isomorphism of stacks $\calM_{n,\calL}\cong\calM_{\GL_{n},\calL}$. 

\subsubsection{Examples} When $G=\SL_{n}$, an $\calL$-twisted $G$-Higgs bundle over $X$ amounts to the same thing as a triple $(\calV,\iota,\varphi)$ where $\calV$ is a vector bundle over $X$ of rank $n$, $\iota:\wedge^{n}\calV\isom\calO_{X}$ and $\varphi: \calV\to\calV\otimes\calL$ satisfies $\Tr(\varphi)=0$.

When $G=\Sp_{2n}$, an $\calL$-twisted $G$-Higgs bundle over $X$ amounts to the same thing as a triple $(\calV,\omega,\varphi)$ where $\calV$ is a vector bundle over $X$ of rank $2n$, $\omega:\wedge^{2}\calV\to\calO_{X}$ is nondegenerate fiberwise,  and $\varphi: \calV\to\calV\otimes\calL$ such that for all local sections $u$ and $v$ of $\calV$,  $\omega(\varphi(u),v)+\omega(u,\varphi(v))=0$ as a local section of $\calL$. 



\subsubsection{Hitchin moduli stack as a cotangent bundle}
In Hitchin's original paper \cite{Hi}, he considered the case where $\calL=\omega_{X}$ is the sheaf of $1$-forms on $X$.   This case is particularly important because $\calM_{G,\omega_{X}}$ is closely related to the cotangent bundle of $\Bun_{G}$. For a point $\calE\in\Bun_{G}(R)$ which is a $G$-torsor over $X_{R}$, the cotangent complex of $\Bun_{G}$ at $\calE$ is given by the derived global sections of the complex $\Ad(\calE)^{\vee}\otimes\omega_{X}$ over $X_{R}$. Using a Killing form on $\frg$ we may identify $\Ad(\calE)^{\vee}$ with $\Ad(\calE)$, therefore the cotangent complex of $\Bun_{G}$ at $\calE$ is $\bR\Gamma(X_{R}, \Ad(\calE)\otimes\omega_{X})$, which lives in degrees $0$ and $1$. In particular, when $\calE\in \Bun_{G}(k)$ has finite automorphism group (e.g., $\calE$ is stable), the {\em Zariski cotangent space} at $\calE$ is $\cohog{0}{X,\Ad(\calE)\otimes\omega_{X}}$, i.e., a cotangent vector of $\Bun_{G}$ at $\calE$ is the same thing as a $\omega_{X}$-twisted Higgs field on $\calE$. Therefore, $T^{*}\Bun_{G}$ (properly defined) and $\calM_{G,\omega_{X}}$ share an open substack $T^{*}\Bun^{s}_{G}$, where $\Bun^{s}_{G}$ is the open substack of stable $G$-bundles.

\subsection{Hitchin fibration}

\subsubsection{Hitchin fibration for $\GL_{n}$}
For an $\calL$-twisted Higgs bundle $(\calV,\varphi)$ on $X$, locally on $X$ we may view $\varphi$ as a matrix with entries which are local sections of $\calL$, and we may take the characteristic polynomial of this matrix. The coefficients of this polynomial are well-defined global sections of $\calL^{\otimes i}$, $1\leq i\leq n$. More intrinsically, $\varphi$ induces a map $\wedge^{i}\varphi:\wedge^{i}\calV\to\wedge^{i}\calV\otimes\calL^{\otimes i}$, and we may take 
\begin{equation*}
a_{i}(\varphi):=\Tr(\wedge^{i}\varphi)\in\cohog{0}{X,\calL^{\otimes i}}.
\end{equation*}
This way we have defined a morphism
\begin{equation*}
f: \calM_{n,\calL}\to\calA_{n,\calL}:=\prod_{i=1}^{n}\cohog{0}{X,\calL^{i}}
\end{equation*}
sending $(\calV,\varphi)$ to $(a_{1}(\varphi),\cdots, a_{n}(\varphi))$. We view $\calA_{n,\calL}$ as an affine space over $k$. The morphism $f$ is called the {\em Hitchin fibration} in the case $G=\GL_{n}$.

\subsubsection{Hitchin fibration in general}\label{sss:Hit fib general} For general connected reductive $G$ as in \S\ref{ss:notation}, the coefficients of the characteristic polynomial in the case of $\GL_{n}$ are replaced with the fundamental $G$-invariant polynomials on $\frg$. Recall that $\frc=\frg\sslash G=\Spec\Sym(\frg^{*})^{G}$. Chevalley's theorem says that $\Sym(\frg^{*})^{G}\cong\Sym(\frt^{*})^{W}$, and the latter is a polynomial ring in $r$ variables. We fix homogeneous generators $f_{1},\cdots, f_{r}$ of $\Sym(\frg^{*})^{G}$ as a $k$-algebra, whose degrees  $d_{1}\leq \cdots \leq d_{r}$ are canonically defined although $f_{i}$ are not canonical.  When $G$ is almost simple, the numbers $e_{i}=d_{i}-1$ are the {\em exponents} of $G$. Viewing $f_{i}$ as a symmetric multilinear function $\frg^{\otimes d_{i}}\to k$ invariant under $G$, for any $G$-torsor $\calE$ over $X$, $f_{i}$ induces a map of the associated bundles
\begin{equation*}
f_{i}: \Ad(\calE)^{\otimes d_{i}}\to\calO_{X}.
\end{equation*}
This further induces
\begin{equation*}
f^{\calL}_{i}:(\Ad(\calE)\otimes\calL)^{\otimes d_{i}}\to\calL^{\otimes d_{i}}.
\end{equation*}
If $\varphi$ is an $\calL$-twisted Higgs field on $\calE$,  we may evaluate $f^{\calL}_{i}$ on the section $\varphi^{\otimes d_{i}}$ of $(\Ad(\calE)\otimes\calL)^{\otimes d_{i}}$ to get
\begin{equation*}
a_{i}(\varphi):=f^{\calL}_{i}(\varphi^{\otimes d_{i}})\in\cohog{0}{X,\calL^{\otimes d_{i}}},  \quad i=1,2,\cdots, r.
\end{equation*}
The assignment $(\calE,\varphi)\mapsto(a_{i}(\varphi))_{1\leq i\leq r}$ defines the {\em Hitchin fibration} for $\calM_{G,\calL}$
\begin{equation*}
f=f_{G,\calL}: \calM_{G,\calL}\to \calA_{G,\calL}:=\prod_{i=1}^{r}\cohog{0}{X,\calL^{\otimes d_{i}}}.
\end{equation*}
The target $\calA_{G,\calL}$ is again viewed as an affine space over $k$, and is called the {\em  Hitchin base}.

A more intrinsic way to define the Hitchin base $\calA_{G,\calL}$ is the following. The affine scheme $\frc=\Spec \Sym(\frg^{*})^{G}$ is  equipped with a $\Gm$-action inducing the grading on its coordinate ring. Let $\Tot^{\times}(\calL)\to X$ be the complement of the zero section in the total space of $\calL$. Consider the $\calL$-twisted version of $\frc$ over $X$:
\begin{equation*}
\frc_{X,\calL}:=(\frc\times \Tot^{\times}(\calL))/\Gm
\end{equation*}
where $\l\in\Gm$ acts by $\l: (c,\wt{x})\mapsto (\l c, \l^{-1}\wt{x})$ on the two coordinates. This is an affine space bundle over $X$ whose fibers are isomorphic to $\frc$. Then $\calA_{G,\calL}$ can be canonically identified with the moduli space of sections of the map $\frc_{X,\calL}\to X$. In particular, every point $a\in\calA_{G,\calL}$ gives a map $a:X\to [\frc/\Gm]$.



\subsubsection{The generically regular semisimple locus} Trivializing $\calL$ at the generic point $\eta$ of $X$ and restricting $a_{i}$ to $\eta$, we have a polynomial $P_{a}(y)=y^{n}-a_{1}y^{n-1}+a_{2}y^{n-2}+\cdots +(-1)^{n}a_{n}\in F[y]$, where $F=k(X)$ is the function field of $X$. When $P_{a}(y)$ is a separable polynomial in $F[y]$, we call such an $a$ {\em generically regular semisimple}. The generic regular semisimplicity of $a$ is equivalent to the nonvanishing of the discriminant $\Delta(P_{a})\in\cohog{0}{X,\calL^{n(n-1)}}$ and therefore it defines an open subscheme $\Ah_{n,\calL}\subset \calA_{n,\calL}$. 

For general $G$, viewing $a\in\calA_{G,\calL}$ as a map $a:X\to [\frc/\Gm]$ (see the end of \S\ref{sss:Hit fib general}), we call $a$ {\em generically regular semisimple} if $a$ sends the generic point of $X$ into the open substack $[\frc^{\rs}/\Gm]$. This defines an open subscheme $\Ah_{G,\calL}\subset\calA_{G,\calL}$ generalizing the construction of $\Ah_{n,\calL}$ above.

\subsubsection{Geometric properties}\label{sss:geom Hit} When $\deg \calL>2g-2$, the stack $\calM|_{\Ah_{G,\calL}}$  is smooth, see \cite[Th 4.14.1]{NgoFL}. In this situation, the morphism $f_{G,\calL}$ is flat over $\Ah_{G,\calL}$, see \cite[Cor 4.16.4]{NgoFL}. When $G$ is semisimple, there is a further open dense subset $\calA^{\textup{ani}}_{G,\calL}\subset \Ah_{G,\calL}$ over which $\calM$ is a Deligne-Mumford stack and the map $f_{G,\calL}$ is proper, see \cite[Prop 6.1.3]{NgoFL}. Comparing to the infinite-dimensionality involved in the geometry of affine Springer fibers, the Hitchin fibration has much nicer geometric properties, and yet it is closely related to the affine Springer fibers, as we shall see in \S\ref{ss:rel Hit ASF}. 


\subsubsection{Generalization}\label{sss:general Hitchin} Let $H$ be a reductive group over $k$ that fits into an exact sequence of reductive groups
\begin{equation*}
1\to H_{1}\to H\to A\to 1.
\end{equation*}
Let $(V,\rho)$ be a representation of $H$. Then we may consider pairs $(\calE,\psi)$ where $\calE$ is an $H$-torsor over $X$ and $\psi$ is a section of the associated bundle $\rho(\calE)$. Alternatively, such a pair is the same as a morphism $X\to [V/H]$. One can prove that there is an algebraic stack $\calM_{H, \rho}$ classifying such pairs. Every $H$-torsor induces an $A$-torsor, hence we have a morphism $\alpha: \calM_{H, \rho}\to \Bun_{A}$. Fix an $A$-torsor $\calE_{A}$ over $X$, we denote the preimage $\alpha^{-1}(\calE_{A})$ by $\calM_{H,\rho,\calE_{A}}$.


To recover the Hitchin moduli stack, we consider the case $H=G\times\Gm$ with $H_{1}=G$ and $A=\Gm$. Let $V=\frg$ with the action $\rho$ of $H$ defined as follows: $G$ acts by the adjoint representation and $\Gm$ acts by scaling on $V$. An $H$-torsor is a pair consisting of a $G$-torsor and a line bundle $\calL$ on $X$. Fixing the line bundle $\calL$ (which is equivalent to fixing an $A=\Gm$-torsor), we get an isomorphism
\begin{equation*}
\calM_{H,\rho, \calL}\cong\calM_{G, \calL}.
\end{equation*}

For general $(1\to H_{1}\to H\to A\to 1,V,\rho,\calE_{A})$ as above, we may define the analog of the Hitchin base as follows. Let $\frc_{H,\rho}=V\sslash H_{1}$. This is the analog of $\frc$, and it carries an action of $A$. Then we form the twisted version of $\frc_{H,\rho}$ over $X$
\begin{equation*}
\frc_{H,\rho,\calE_{A}}:=\calE_{A}\twtimes{A}\frc_{H,\rho}
\end{equation*}
Then we define $\calA_{H,\rho,\calE_{A}}$ to be the moduli space of sections to the map $\frc_{H,\rho,\calE_{A}}\to X$. The morphism $[V/H]\to [\frc_{H,\rho}/A]$ then induces the analog of the Hitchin fibration
\begin{equation*}
f_{H,\rho,\calE_{A}}: \calM_{H,\rho,\calE_{A}}\to \calA_{H,\rho,\calE_{A}}
\end{equation*}


\subsubsection{Example}\label{sss:MHilb} Consider $H=\GL(U)\times\Gm\times\Gm$, and $V=\End(U)\oplus U^{*}$. The action $\rho(g,s_{1}, s_{2})$ on $V$ is given by $(A,u^{*})\mapsto (s_{1}gAg^{-1}, s_{2}gu^{*})$. The moduli stack $\calM_{H,\rho}$ then maps to $\Pic(X)\times\Pic(X)$ by remembering only the two $\Gm$-torsors. Fixing $(\calL_{1},\calL_{2})\in \Pic(X)\times\Pic(X)$, its  preimage  in $\calM_{H,\rho}$ classifies triples $(\calU,\varphi,\beta)$ where $(\calU,\varphi)$ is a Higgs bundle over $X$ of rank $n=\dim V$, $\beta$ is an $\calO_{X}$-linear map $\calU\to \calL_{2}$. The Hitchin base in this case is the same as the classical Hitchin base $\calA_{n,\calL_{1}}$. Later in \S\ref{sss:Hilb} we will relate this moduli space to Hilbert schemes of curves.


